% Generated from main.typ. Typst is authoritative; regeneration overwrites this file.
\documentclass{qlnotes-export}

\title{MATH 597: Measure Theory}
\subtitle{Typst-first course notes and worked homeworks}
\author{Qiulin Fan}
\date{Winter 2025}

\begin{document}
\mainmatter
\chapter{sigma-algebra 与 measure}\label{sigma-algebra-ux4e0e-measure}

\section{\texorpdfstring{\(\sigma\)-algebra {[}Fol 1.2{]}}{\textbackslash sigma-algebra {[}Fol 1.2{]}}}

我们 (见 my Math 395 notes) 已经证明: 在 \(\mathbb{R}\) 上不存在一个 measure function \(\mu:\mathcal{P}({\mathbb{R}})\rightarrow\lbrack 0,\infty\rbrack\) satisfying:

\begin{enumerate}
\item
  \(\mu(\varnothing) = 0\);
\item
  translate invariant
\item
  countably additivite
\end{enumerate}

因而, 对于比如 \(\mathbb{R}\) 的这种无法在其幂集上定义良好的 measure function 的集合, 我们要定义一个 \(\mathcal{A} \subseteq \mathcal{P}(X)\), 使得我们能在这个 power set 的子集上, 定义一个 make sense 的 measure.\\
\strut \\
首先, 为了对于一个任意的集合 \(X\) 都能在其上定义 measure, 我们要考虑在 \(X\) 的一个什么样的子集簇上有希望定义这样的 measure.

\begin{definition}{\protect\phantomsection\label{kn-sigma-algebra}{\textbf{\(\sigma\)-algebra}}}

对于 set \(X\), \(S \subseteq \mathcal{P}(X)\) 被称为 \(X\) 上的一个 \(\sigma\)-algebra, if 其满足:

\begin{enumerate}
\item
  \(\varnothing \in X\);
\item
  \textbf{closed under complement}: if \(E \in S\) then \(X\backslash E \in S\);
\item
  \textbf{closed under countable union}: if \(E_{1},E_{2},\cdots \in S\) then \(\bigcup_{k = 1}^{\infty}E_{k} \in S\).
\end{enumerate}

如果第三条并不满足, 而是只满足 \textbf{closed under finite union}, 则称 \(S\) 是 \(X\) 上的一个 \protect\phantomsection\label{kn-algebra-of-sets}{\textbf{algebra of sets}} . 当然, \(\sigma\)-algebra 是比 algebra 严格更强的条件.

\end{definition}

我们定义 \(X\) 的一个子集簇为一个 \(\sigma\)-algebra 如果它包含空集并 closed under complement and countable union. 但这并不是 \(\sigma\)-algebra 的全部性质. 这三个性质还蕴涵了: \(\sigma\)-algebra 也一定包含 \(X\), 且 \textbf{closed under set difference}, \textbf{symmetric difference} 以及 \textbf{countable intersection}.\\
\textbf{对于 algebra, 它也有以上的所有性质的 finite version.}

\begin{theorem}{\protect\phantomsection\label{kn-sigma-algebra-closure-properties}{\textbf{\(\sigma\)-algebra
closure properties: set difference, symmetric difference and countable
intersection}}}

Let \(S\) be a \(\sigma\)-algebra on set \(X\).\\
Claim:

\begin{enumerate}
\item
  \(X \in S\)

  \begin{proof}

  Directly from def.

  \end{proof}
\item
  \(D,E \in S\Longrightarrow D \cup E,D \cap E,D\backslash E \in S\)

  \begin{proof}

  union: from def by leaving others as \(\varnothing\);\\
  intersection:

  \[(D \cap E)^{C} = D^{C} \cup E^{C} \in S\]

  setminus:

  \[D\backslash E = D \cap (X\backslash E) \in S\]

  \end{proof}
\item
  \(D,E \in S\Longrightarrow D\Delta S \in S\)

  \begin{proof}

  \[D\Delta E = (D\backslash E)\bigcup(E\backslash D)\]

  \end{proof}
\item
  \(A_{1},A_{2},\cdots \in S\Longrightarrow\bigcap_{i = 1}^{\infty}A_{i} \in S\)

  \begin{proof}

  \[(\bigcap\limits_{n = 1}^{\infty})^{C} = \bigcup\limits_{n = 1}^{\infty}E_{n}^{C} \in S\]

  \end{proof}
\end{enumerate}

\end{theorem}

\begin{qlremark}{Remark}

我们发现 \(\sigma\)-algerbra 很像是 topology. 实际上 \(\sigma\)-algerbra 和 topology 的区别就是: \(\sigma\)-algebra 只保证了 closed under countable union 而 topology closed under any union; topology 只保证 closed under finite intersection 而 \(\sigma\)-algebra closed under countable intersection.

\end{qlremark}

\begin{lemma}{\protect\phantomsection\label{kn-sigma-algebra-intersection-sigma-algebra}{\textbf{任意
\(\sigma\)-algebra 的 intersection 仍是 \(\sigma\)-algebra}}}

Let \(\left\{ S_{\alpha} \right\}_{\alpha \in A}\) be a collection of \(\sigma\)-algebra on \(X\), then \(\bigcap_{\alpha \in A}S_{\alpha}\) is a \(\sigma\)-algebra on \(X\).

\end{lemma}

\begin{proof}

这是个 trivial proof. 但是它具有一定理解上的启发.\\
我们对 \(\sigma\)-algebra 有一个直观理解: 如果我们想把一些集合做成一个 \(\sigma\)-algebra, 那么首先我们把它们的补集放进这个 \(\sigma\)-algebra 里, 其次我们把这些集合的 up to countable 的任意组合的并集也放进这个 \(\sigma\)-algebra 里.\\
因而即便我们把一些 \(\sigma\)-algebra 给 intersect 起来, 其中每个集合的补集和这些集合的 up to ctbl 的任意组合的并集也在这个 intersection 里.\\
这是个重要的直观理解. 我们想到, 如果我们要把一个 sigma-algebra 里的一部分去掉，并保持它仍然是一个 sigma-algebra，那么我们得把这些集合的补集, 以及能够 ctbly union 成这些集合的小集合也去掉, 并对这些小集合也 recursively 进行这个操作.

\end{proof}

\begin{corollary}{\protect\phantomsection\label{kn-unique-smallest-sigma-algebra-containing-a-collection-of-subsets}{\textbf{unique
smallest \(\sigma\)-algebra containing a collection of subsets}}}

Given \(\varepsilon \subseteq \mathcal{P}(X)\)

\[< \varepsilon > := \bigcap\limits_{\varepsilon \subseteq S \subseteq \mathcal{P}(X),\, S\ \text{is}\ \sigma\ \text{-algebra on}\ X}S\]

\end{corollary}

\begin{definition}{\protect\phantomsection\label{kn-sigma-algebra-generated-by-a-subset}{\textbf{\(\sigma\)-algebra
generated by a subset}}}

We call

\[< \varepsilon > := \bigcap\limits_{\varepsilon \subseteq S \subseteq \mathcal{P}(X),\, S\ \text{is}\ \sigma\ \text{-algebra on}\ X}S\]

\textbf{the \(\sigma\)-algebra generated by \(\varepsilon\)}

\end{definition}

\section{\texorpdfstring{Borel \(\sigma\)-algebra on \(\mathbb{R}\) and measure {[}Fol 1.2, finished; 1.3{]}}{Borel \textbackslash sigma-algebra on \textbackslash mathbb\{R\} and measure {[}Fol 1.2, finished; 1.3{]}}}

Recall: the \(\sigma\)-algebra generated by \(\varepsilon\)

\[< \varepsilon > := \bigcap\limits_{\varepsilon \subseteq S \subseteq \mathcal{P}(X),\, S\ \text{is}\ \sigma\ \text{-algebra on}\ X}S\]

is the smallsest \(\sigma\)-algebra containing \(\varepsilon\).

\begin{example}

\[< \left\{ E \right\} > = \left\{ {\varnothing,E,E^{c},X} \right\}\]

\end{example}

% qlnotes: kind=lemma; id=lem-01-sigma-algebra-measure-inclusion-properties-of-generated-sigma-algebra
\begin{lemma}{\protect\phantomsection\label{kn-inclusion-properties-of-generated-sigma-algebra}{\textbf{inclusion
properties of generated \(\sigma\)-algebra}}}\label{lem-01-sigma-algebra-measure-inclusion-properties-of-generated-sigma-algebra}

\begin{enumerate}
\item
  if \(\mathcal{E} \subseteq \mathcal{A}\) where \(\mathcal{A}\) is a \(\sigma\)-algebra, then \(< \mathcal{E} > \subseteq \mathcal{A}\).
\item
  if \(\mathcal{E} \subseteq \mathcal{F}\), then \(< \mathcal{E} > \subseteq < \mathcal{F} >\).
\item
  if \(\mathcal{E} \subseteq < \mathcal{F} >\), then \(< \mathcal{E} > \subseteq < \mathcal{F} >\).
\end{enumerate}

\end{lemma}

\begin{proof}

trivial.

\end{proof}

\begin{definition}{\protect\phantomsection\label{kn-borel-sigma-algebra-defined-on-a-topological-space}{\textbf{Borel
\(\sigma\)-algebra defined on a topological space}}}

For topological space \((X,\mathcal{T})\), we define:

\[\mathcal{B}_{X} := < \mathcal{T} >\]

\end{definition}

\textbf{Borel \(\sigma\)-algebra} on a topological space 就是 \(\sigma\)-algebra generated by the topology. Its members are called \textbf{Borel sets}. 当然, 所有的 open sets 和 closed sets 都是 Borel sets.

\subsection{\texorpdfstring{generating Borel \(\sigma\)-algebra on \(\mathbb{R}\)}{generating Borel \textbackslash sigma-algebra on \textbackslash mathbb\{R\}}}\label{generating-borel-sigma-algebra-on-mathbbr}

\begin{example}

Let \(\mathcal{E}_{1}\): \(\mathbb{R}\) 上所有的 open intervals;\\
\(\mathcal{E}_{2}\): \(\mathbb{R}\) 上所有的 closed intervals;\\
\(\mathcal{E}_{3}\): \(\mathbb{R}\) 上所有的左开右闭 intervals;\\
\(\mathcal{E}_{4}\): \(\mathbb{R}\) 上所有的左闭右开 intervals;\\
\(\mathcal{E}_{5}\): \(\mathbb{R}\) 上所有的左开右无界 intervals;\\
\(\mathcal{E}_{6}\): \(\mathbb{R}\) 上所有的左闭右无界 intervals;\\
\(\mathcal{E}_{7}\): \(\mathbb{R}\) 上所有的左无界右开 intervals;\\
\(\mathcal{E}_{8}\): \(\mathbb{R}\) 上所有的左无界右闭 intervals;\\
\(\bigcup_{i = 1,\cdots,8}\mathcal{E}_{i}\) 即 \(\mathbb{R}\) 上的所有形式的 interals.

\begin{lemma}

任意以上 \(\mathcal{E}_{i},i = 1,\cdots,8\) 都可以 generate \(\mathcal{B}_{\mathbb{R}}\)

\end{lemma}

\end{example}

\begin{proof}

我们 recall: 所有的 countable 以及 second countable 的 topological space 都具有 \textbf{Lindelöf property}: 任意 open covering 都存在一个 countable 的 subcovering.\\
Lindelöf property 的一个推论就是, 在具有 Lindelöf property 的 metric space 或者 second countable 的 space 中, 任意 open set 都可以写成 countable 个 open balls 的 union.\\
我们在 elementary 的 real analysis 中已经学过, \(\lbrack a,b) = \cap_{n \geq 1}(a - 1/n,b)\), 以其作为例子, 这些 intervals 彼此之间都可以相互转换.

\end{proof}

\subsection{measure}\label{measure}

\begin{definition}{\protect\phantomsection\label{kn-measurable-space-and-measure-space}{\textbf{measurable
space}} and
\protect\phantomsection\label{kn-measure-space}{\textbf{measure space}}}

Let \(X\) be a set, \(\mathcal{M}\) be a \(\sigma\)-algebra on \(X\). We call \((X,\mathcal{M})\) a measurable space.\\
A \protect\phantomsection\label{kn-measure}{\textbf{measure}} on this measurable space is a function \(\mu:\mathcal{M}\rightarrow\lbrack 0,\infty)\) satisfying:

\begin{enumerate}
\item
  \(\mu(\varnothing) = 0\)
\item
  countable additive:

  \[\mu(\bigcup\limits_{i = 1}^{\infty}E_{i}) = \sum\limits_{i = 1}^{\infty}\mu(E_{i})\]

  for disjoint seq of \(E_{i} \in \mathcal{M}\).
\end{enumerate}

我们称 \((X,\mathcal{M},\mu)\) 为一个 measure space.

\end{definition}

\begin{qlremark}{Remark}

一个 probability space 就是一个 measure space, satisfying \(\mu(X) = 1\).

\end{qlremark}

\begin{example}

\begin{enumerate}
\item
  对于任意的 \((X,\mathcal{M})\), 我们可以定义:

  \[\mu(A) := \# A( \in {\mathbb{Z}}_{\geq 0} \cup \left\{ \infty \right\})\]

  这个 measure 叫做 \textbf{counting measure}.
\item
  Fix \(x_{0} \in M\), 可以 define

  \[\begin{matrix}
  {\mu(A) := \delta_{x} := \{ 1\text{, if}\ x_{0} \in A} \\
  {0\text{, if}\ x_{0} \notin A}
  \end{matrix}\]

  这个 measure 叫做 the \textbf{Dirac measure at \(x_{0}\)}.
\item
  给定一个 \(X\) 上的函数 \(f:X\rightarrow\lbrack 0,\infty)\), 我们可以通过这个函数来定义:

  \[\mu(A) := \sum\limits_{x \in A}f(x)\]

  这个测度依赖于函数值来表示每个点的单点集的 measure, 并通过一个集合上所有点的单点集 measure 相加得到这个集合在这个函数下的 measure. (缺点: 我们已经知道, 如果一个函数在一个集合上的正集是 uncountable 的, 那么这个集合上的这个测度一定是 \(\infty\).)
\end{enumerate}

\end{example}

以下是 measure function 由它的定义的两条性质(空集为0以及 ctbl additivity)推导出的一些基本性质:

\begin{lemma}{\protect\phantomsection\label{kn-measure-is-finitely-additive}{\textbf{measure
is finitely additive}}}

Measure is finitely additive.

\end{lemma}

\begin{proof}

显然, ctbl additive implies finite additive.

\end{proof}

\begin{qlremark}{Remark}

反向则不成立. 这让我们想起: Jordan measure 和 Lebesgue measure.

\end{qlremark}

\begin{lemma}

\(A,B \in \mathcal{M}\Longrightarrow\)

\[\mu(A) + \mu(B) = \mu(A \cap B) + \mu(A \cup B)\]

\end{lemma}

\begin{proof}

\[A \cup B = (A\backslash B) \sqcup (A \cap B) \sqcup (B\backslash A)\]

而后使用 finite additive 可得. 这是一个 direct corollary of countable additivity.

\end{proof}

\begin{corollary}

\(A,B \in \mathcal{M},A \subseteq B,\mu(A) < \infty\Longrightarrow\)

\[\mu(B\backslash A) = \mu(B) - \mu(A)\]

\end{corollary}

\begin{theorem}{\protect\phantomsection\label{kn-properties-of-measure}{\textbf{properties
of measure}}}

对于任何 measure space \((X,\mathcal{M},\mu)\):

\begin{enumerate}
\item
  \textbf{monotonicity}: \(A \subseteq B \in \mathcal{M}\Longrightarrow\mu(A) \leq \,\mu(B)\)

  \begin{proof}

  trivial.

  \end{proof}
\item
  \textbf{countable subadditivity}:

  \[\mu(\bigcup\limits_{i = 1}^{\infty}A_{i}) \leq \sum\limits_{i = 1}^{\infty}\mu(A_{i})\]

  \begin{proof}

  By setting \(B_{i} = A_{i}\backslash\bigcup_{j = 1}^{i - 1}A_{j}\), 而后通过 ctbl disjoint additivity 与 monotonicity 可得

  \end{proof}
\item
  \textbf{continuous from above}: 如果 \(A_{i} \subseteq A_{i + 1}\forall i \geq 2\Longrightarrow\)

  \[\mu(\bigcup\limits_{i = 1}^{\infty}A_{i}) = \lim\limits_{i\rightarrow\infty}\mu(A_{i})\]

  \begin{proof}

  使用 same trick as 2.

  \end{proof}
\item
  \textbf{countinuous from below}: 如果 \(A_{i} \supseteq A_{i + 1}\forall i\) 且存在某个 \(j\) 使得 \(\mu(A_{i}) < \infty\), 则

  \[\mu(\bigcap\limits_{i = 1}^{\infty}A_{i}) = \lim\limits_{n\rightarrow\infty}\mu(A_{n})\]

  \begin{proof}

  前面的都无视, 直到第一个 measure \(< \infty\) 的集合, 是可能出现在最后的 intersection 里的最大集合. 我们 Fix 这个 \(A_{j}\). 通过构造补集的方式, 把交转为并, 从而用 (3) 得证. Define: \(E_{i} := A_{j}\backslash A_{i}\forall i \geq j\) 从而

  \[\bigcup\limits_{i = j}^{\infty}E_{i} = A_{j}\backslash(\bigcap\limits_{i = j}^{\infty}A_{i})\]

  进而

  \[\mu(\bigcup\limits_{i = j}^{\infty}E_{i}) = \mu(A_{j}) - \mu(\bigcap\limits_{i = j}^{\infty}A_{i})\]

  进而 by (3)

  \[\mu(\bigcap\limits_{i = 1}^{\infty}A_{i}) = \mu(\bigcap\limits_{i = j}^{\infty}A_{i}) = \mu(A_{j}) - \lim\limits_{i\rightarrow\infty}\mu(E_{i}) = \mu(A_{j}) - \lim\limits_{i\rightarrow\infty}(\mu(A_{j}) - \mu(A_{i})) = \lim\limits_{i\rightarrow\infty}A_{i}\]

  \end{proof}
\end{enumerate}

\end{theorem}

Recall: the \(\sigma\)-algebra generated by \(\varepsilon\)

\[< \varepsilon > := \bigcap\limits_{\varepsilon \subseteq S \subseteq \mathcal{P}(X),\, S\ \text{is}\ \sigma\ \text{-algebra on}\ X}S\]

is the smallsest \(\sigma\)-algebra containing \(\varepsilon\).

以下是 measure function 由它的定义的两条性质(空集为0以及 ctbl additivity)推导出的一些基本性质:

\begin{lemma}{\href{https://qiulinfan.github.io/qlblog/notes/math/measure-theory/\#kn-measure-is-finitely-additive}{measure
is finitely additive}}

Measure is finitely additive.

\end{lemma}

\begin{proof}

显然, ctbl additive implies finite additive.

\end{proof}

\begin{qlremark}{Remark}

反向则不成立. 这让我们想起: Jordan measure 和 Lebesgue measure.

\end{qlremark}

\begin{lemma}

\(A,B \in \mathcal{M}\Longrightarrow\)

\[\mu(A) + \mu(B) = \mu(A \cap B) + \mu(A \cup B)\]

\end{lemma}

\begin{proof}

\[A \cup B = (A\backslash B) \sqcup (A \cap B) \sqcup (B\backslash A)\]

而后使用 finite additive 可得. 这是一个 direct corollary of countable additivity.

\end{proof}

\begin{corollary}

\(A,B \in \mathcal{M},A \subseteq B,\mu(A) < \infty\Longrightarrow\)

\[\mu(B\backslash A) = \mu(B) - \mu(A)\]

\end{corollary}

\begin{theorem}{\href{https://qiulinfan.github.io/qlblog/notes/math/measure-theory/\#kn-properties-of-measure}{properties
of measure}}

对于任何 measure space \((X,\mathcal{M},\mu)\):

\begin{enumerate}
\item
  \textbf{monotonicity}: \(A \subseteq B \in \mathcal{M}\Longrightarrow\mu(A) \leq \,\mu(B)\)

  \begin{proof}

  trivial.

  \end{proof}
\item
  \textbf{countable subadditivity}:

  \[\mu(\bigcup\limits_{i = 1}^{\infty}A_{i}) \leq \sum\limits_{i = 1}^{\infty}\mu(A_{i})\]

  \begin{proof}

  By setting \(B_{i} = A_{i}\backslash\bigcup_{j = 1}^{i - 1}A_{j}\), 而后通过 ctbl disjoint additivity 与 monotonicity 可得

  \end{proof}
\item
  \textbf{continuous from above}: 如果 \(A_{i} \subseteq A_{i + 1}\forall i \geq 2\Longrightarrow\)

  \[\mu(\bigcup\limits_{i = 1}^{\infty}A_{i}) = \lim\limits_{i\rightarrow\infty}\mu(A_{i})\]

  \begin{proof}

  使用 same trick as 2.

  \end{proof}
\item
  \textbf{countinuous from below}: 如果 \(A_{i} \supseteq A_{i + 1}\forall i\) 且存在某个 \(j\) 使得 \(\mu(A_{i}) < \infty\), 则

  \[\mu(\bigcap\limits_{i = 1}^{\infty}A_{i}) = \lim\limits_{n\rightarrow\infty}\mu(A_{n})\]

  \begin{proof}

  前面的都无视, 直到第一个 measure \(< \infty\) 的集合, 是可能出现在最后的 intersection 里的最大集合. 我们 Fix 这个 \(A_{j}\). 通过构造补集的方式, 把交转为并, 从而用 (3) 得证. Define: \(E_{i} := A_{j}\backslash A_{i}\forall i \geq j\) 从而

  \[\bigcup\limits_{i = j}^{\infty}E_{i} = A_{j}\backslash(\bigcap\limits_{i = j}^{\infty}A_{i})\]

  进而

  \[\mu(\bigcup\limits_{i = j}^{\infty}E_{i}) = \mu(A_{j}) - \mu(\bigcap\limits_{i = j}^{\infty}A_{i})\]

  进而 by (3)

  \[\mu(\bigcap\limits_{i = 1}^{\infty}A_{i}) = \mu(\bigcap\limits_{i = j}^{\infty}A_{i}) = \mu(A_{j}) - \lim\limits_{i\rightarrow\infty}\mu(E_{i}) = \mu(A_{j}) - \lim\limits_{i\rightarrow\infty}(\mu(A_{j}) - \mu(A_{i})) = \lim\limits_{i\rightarrow\infty}A_{i}\]

  \end{proof}
\end{enumerate}

\end{theorem}
\end{document}
